% -*- coding: utf-8 -*-
% !TEX program = xelatex
\documentclass[plain,noanswer,medmath]{../../template/jnuexam}
\usepackage{../../template/kysx}

\begin{document}


\examtitle{1994}{5}


\loadproblems{3}{1994P3}%载入试卷三
\loadproblems{4}{1994P4}%载入试卷四


\makepart{填空题}{本题共5小题，每小题3分，满分15分}


\useproblem{4}{1}{1}%【同试卷四第一(1)题】


\useproblem{4}{1}{2}%【同试卷四第一(2)题】


\useproblem{4}{1}{3}%【同试卷四第一(3)题】


\useproblem{4}{1}{4}%【同试卷四第一(4)题】


\begin{problem}
假设一批产品中一，二，三等品各占$60\%$，$30\%$，$10\%$，
从中随意取出一件，结果不是三等品，则取到的是一等品的概率为 \fillin{}．
\end{problem}


\makepart{选择题}{本题共5小题，每小题3分，满分15分}


\useproblem{3}{2}{4}%【同试卷三第二(4)题】【同试卷四第二(1)题】


\begin{problem}
设函数$f(x)$在闭区间$[a,b]$上连续，且$f(x) > 0$，
则方程$\int_a^xf(t)\dt + \int_b^x\frac{1}{f(t)}\dt = 0$
在开区间$(a,b)$内的根有\pickout{}
\begin{abcd}
\item $0$个．      
\item $1$个．      
\item $2$个．      
\item 无穷多个．
\end{abcd}
\end{problem}


\begin{problem}
设$A$, $B$都是$n$阶非零矩阵，且$AB = O$，则$A$和$B$的秩\pickout{}
\begin{abcd}
\item 必有一个等于零．					
\item 都小于$n$．
\item 一个小于$n$，一个等于$n$．			
\item 都等于$n$．
\end{abcd}
\end{problem}


\begin{problem}
设有向量组$\alpha_1 = (1, - 1,2,4)$, $\alpha_2 = (0,3,1,2)$, $\alpha_3 = (3,0,7,14)$,
$\alpha_4 = (1, - 2,2,0)$, $\alpha_5 = (2,1,5,10)$，则该向量组的极大线性无关组是\pickout{}
\begin{abcd}
\item $\alpha_1$, $\alpha_2$, $\alpha_3$．
\item $\alpha_1$, $\alpha_2$, $\alpha_4$．
\item $\alpha_1$, $\alpha_2$, $\alpha_5$．
\item $\alpha_1$, $\alpha_2$, $\alpha_4$, $\alpha_5$．
\end{abcd}
\end{problem}


\useproblem{4}{2}{4}%【同试卷四第二(4)题】


\makepart{}{本题满分5分}

\begin{problem}
求极限$\lim_{x \to \infty} \left[x - x^2\ln \left(1 + \frac{1}{x} \right) \right]$．
\end{problem}


\makepart{}{本题满分5分}%【同试卷四第五题】

\useproblem{4}{5}{0}%


\makepart{}{本题满分6分}

\begin{problem}
设$\frac{\sin x}{x}$是$f(x)$的一个原函数，求$\int x^3 f'(x)\dx$．
\end{problem}


\makepart{}{本题满分8分}

\begin{problem}
某养殖场饲养两种鱼，若甲种鱼放养$x$（万尾），乙种鱼放养$y$（万尾），
收获时两种鱼的收获量分别为$(3 - \alpha x - \beta y)x$和$(4 - \beta x - 2\alpha y)y$（$\alpha > \beta > 0$)．
求使产鱼总量最大的放养数．
\end{problem}


\makepart{}{本题满分8分}%【类似试卷四第七题】

\begin{problem}
已知曲线$y = a\sqrt{x}$（$a > 0$）与曲线$y = \ln \sqrt{x}$在点$(x_0,y_0)$处有公共切线，求：
\begin{enumerate}
\item 常数$a$及切点$(x_0,y_0)$；
\item 两曲线与$x$轴围成的平面图形的面积$S$．
\end{enumerate}
\end{problem}


\makepart{}{本题满分7分}%【类似试卷四第六题】

\useproblem{4}{6}{0}%

%\begin{problem}
%设函数$f(x)$可导，且$f(0) = 0$, $F(x) = \int_0^x t^{n - 1} f(x^n - t^n)\dt$，
%证明$\lim_{x \to 0} \frac{F(x)}{x^{2n}} = \frac{1}{2n}f'(0)$．
%\end{problem}


\makepart{}{本题满分8分}

\begin{problem}
设$\alpha_1$, $\alpha_2$, $\alpha_3$是齐次线性方程组$Ax = 0$的一个基础解系，
证明：$\alpha_1 + \alpha_2$, $\alpha_2 + \alpha_3$, $\alpha_3 + \alpha_1$也是该方程组的一个基础解系．
\end{problem}


\makepart{}{本题满分8分}%【同试卷四第十题】

\useproblem{4}{10}{0}%


\makepart{}{本题满分7分}

\begin{problem}
假设随机变量$X$的概率密度为$f(x) = \begin{cases}
  2x, & 0 < x < 1, \\
   0, & \text{其他}.
\end{cases}$现在对$X$进行$n$次独立重复观测，
以$V_n$表示观测值不大于$0.1$的次数，试求随机变量$V_n$的概率分布．
\end{problem}


\makepart{}{本题满分8分}%【同试卷四第十二题】

\useproblem{4}{12}{0}%


\end{document}
